\kchapter{Troubleshooting Playbook}{ep}

When something is wrong, resist the urge to guess. Work from the outside in,
and let Kubernetes tell you what it knows. Nearly every problem is solved by
some combination of \kc{get}, \kc{describe}, \kc{logs}, and \kc{events}.

\section{The universal first five commands}

\begin{bashcode}{terminal}
kubectl get pods -n <ns> -o wide                 # 1. what state is everything in?
kubectl describe pod <pod> -n <ns>               # 2. read the Events at the bottom
kubectl logs <pod> -n <ns> [--previous]          # 3. what did the app say?
kubectl get events -n <ns> --sort-by=.lastTimestamp   # 4. what did the cluster do?
kubectl get endpointslices -n <ns> -l kubernetes.io/service-name=<svc>   # 5. is the Service wired up?
\end{bashcode}

\section{Decision tree: my app is not reachable}

\begin{center}
\resizebox{\linewidth}{!}{%
\begin{tikzpicture}[
  q/.style={draw=k8sblue, fill=k8slight, rounded corners=5pt, align=center,
            font=\scriptsize\bfseries, text=k8snavy, minimum width=3.4cm, minimum height=1cm},
  a/.style={fill=k8sorange!15, draw=k8sorange, rounded corners=5pt, align=center,
            font=\scriptsize, text width=3.6cm, minimum height=1cm},
  >={Stealth[length=2.4mm]}]
  \node[q] (q1) at (0,0)      {Are the Pods Running\\and READY 1/1?};
  \node[q] (q2) at (0,-2.2)   {Does the Service have\\endpoints?};
  \node[q] (q3) at (0,-4.4)   {Does \kc{curl svc:port}\\work from inside?};
  \node[q] (q4) at (0,-6.6)   {Does the Gateway/Ingress\\route show Accepted?};
  \node[a] (a1) at (5.6,0)    {Fix the Pod first: see the status table below.};
  \node[a] (a2) at (5.6,-2.2) {Selector does not match Pod labels, or Pods are not Ready.};
  \node[a] (a3) at (5.6,-4.4) {Wrong \kc{targetPort}, app listening on 127.0.0.1, or a NetworkPolicy blocks it.};
  \node[a] (a4) at (5.6,-6.6) {Check \kc{kubectl describe httproute}: parentRef, hostname, backend name/port.};
  \node[a, fill=k8sgreen!15, draw=k8sgreen] (ok) at (0,-8.8) {Check DNS, external LB, firewall, and TLS certificates.};
  \foreach \x/\y in {q1/q2, q2/q3, q3/q4, q4/ok} \draw[->, thick, k8sgreen] (\x) -- node[left, font=\tiny]{yes} (\y);
  \foreach \x/\y in {q1/a1, q2/a2, q3/a3, q4/a4} \draw[->, thick, k8sred] (\x) -- node[above, font=\tiny]{no} (\y);
\end{tikzpicture}}
\end{center}

\section{Pod status reference}

\begin{center}
\ktablestart
\begin{tabularx}{\linewidth}{P{3.4cm} L L}
\kheader \kth{Symptom} & \kth{Likely cause} & \kth{What to check / do} \\
\kc{Pending} & Not enough CPU/memory; nodeSelector/affinity/taint mismatch; PVC unbound. &
   \kc{describe pod} events (``0/3 nodes are available\ldots''); \kc{get pvc}; lower requests or add nodes. \\
\kc{ImagePullBackOff} \kc{ErrImagePull} & Typo in image/tag; private registry; rate limits. &
   \kc{describe pod}; verify the tag exists; add \kc{imagePullSecrets}; for kind, \kc{kind load}. \\
\kc{CrashLoopBackOff} & App exits on start: bad config, missing env, failing dependency, wrong command. &
   \kc{logs --previous}; check env/ConfigMap/Secret; run the image locally. \\
\kc{OOMKilled} & Memory limit too low or a leak. &
   \kc{describe pod} (Last State); raise the limit; profile memory. \\
\kc{CreateContainerConfigError} & Missing ConfigMap/Secret or key. &
   \kc{describe pod}; \kc{get cm,secret}; check names after Kustomize hashing. \\
Running but 0/1 READY & Readiness probe failing. &
   \kc{describe pod} probe messages; hit the endpoint with \kc{exec}/\kc{port-forward}. \\
Restarts keep rising & Liveness probe too strict, or real crashes. &
   Add/relax a \kc{startupProbe}; check \kc{logs --previous}. \\
\kc{Terminating} forever & Finalizers, or an unreachable node. &
   \kc{get pod -o yaml} (finalizers); as a last resort \kc{delete --grace-period=0 --force}. \\
\kc{Evicted} & Node under memory/disk pressure. &
   \kc{describe node}; set requests; clean disk; check ephemeral storage. \\
\end{tabularx}
\end{center}

\section{Networking checks}

\begin{bashcode}{terminal}
# A throw-away debug Pod with networking tools
kubectl run netshoot --rm -it --image=nicolaka/netshoot -- bash
#   nslookup api.taskflow-dev.svc.cluster.local
#   curl -v http://api.taskflow-dev/api/info
#   nc -zv postgres.taskflow-dev 5432

kubectl get networkpolicy -n <ns>
kubectl -n kube-system logs -l k8s-app=kube-dns          # CoreDNS problems
kubectl get gateway,httproute -A                          # route status and conditions
\end{bashcode}

\begin{warningbox}
If the namespace enforces the \kc{restricted} Pod Security profile, a plain
\kc{kubectl run} debug Pod will be rejected. Debug from another namespace, use
\kc{kubectl debug} on an existing Pod, or pass a compliant \kc{securityContext}
with \kc{--overrides}.
\end{warningbox}

\section{Node and cluster checks}

\begin{bashcode}{terminal}
kubectl get nodes                                  # NotReady?
kubectl describe node <node>                       # Conditions: MemoryPressure, DiskPressure, PIDPressure
kubectl get --raw='/readyz?verbose'                # API server health checks
kubectl get pods -n kube-system                    # control-plane and add-on Pods healthy?
\end{bashcode}

\section{Habits that prevent incidents}

\begin{itemize}
  \item Always \kc{kubectl diff} before \kc{apply} in shared environments.
  \item Show the current context in your shell prompt.
  \item Never use \kc{latest} tags; pin versions or digests.
  \item Set requests and limits on everything.
  \item Keep everything in Git; avoid \kc{kubectl edit} in production.
  \item Write a runbook entry for every incident you solve.
\end{itemize}
